\subsubsection{アーキテクチャの実践、方法、戦術}
アーキテクチャの実践には、多くの方法がある。方法とは、関心事の発見、要求整理、候補解作成、文書化、レビュー、評価をどの順序で行うかを定める作業手順である。戦術とは、特定の品質特性を満たすための小さな設計上の手段である。

たとえば、可用性を高めるには冗長化、監視、障害切替が使われることがある。性能を高めるにはキャッシュ、非同期処理、並列化が使われることがある。変更容易性を高めるには抽象化、疎結合、情報隠蔽が使われることがある。セキュリティを高めるには認証、認可、入力検証、監査ログが使われることがある。

\begin{statementbox}{注意}
戦術は魔法の解決策ではない。どの戦術にも副作用がある。キャッシュは性能を高めるが、整合性管理を難しくすることがある。冗長化は可用性を高めるが、費用と運用複雑性を増やすことがある。
\end{statementbox}
